\section{Proof of the sharp bounds for Boolean phase lines}
\label{bound:appendix}

This appendix proves Theorem~\ref{bound:theorem}: tetrahedral regularity fixes the signed local components, coefficient comparison gives the four radius laws, and low-dimensional character identities determine which laws are possible.

\subsection{From tetrahedral regularity to equal signed components}
The local action of the Pauli orbit sends a seed Bloch vector $(x,y,z)$
to the four vectors
\begin{equation}
 (x,y,z),\quad(-x,-y,z),\quad(x,-y,-z),\quad(-x,y,-z).
 \label{bound:orbit-four}
\end{equation}
The label map is onto these four sign patterns.  The three distinct
off-diagonal inner products are
$-x^2-y^2+z^2$, $x^2-y^2-z^2$, and $-x^2+y^2-z^2$.
They are equal precisely when $x^2=y^2=z^2$.
A common radius at all sites then makes the squares of all $3n$
signed components equal.  At the specified base point each signed
component equals $c_n>0$.  By continuity their signs remain positive
on a sufficiently small interval, and the components themselves are
equal there.

For Boolean $d$, every term in a local expectation has a phase
$e^{it[d(z)-d(z')]}$ with $d(z)-d(z')\in\{0,\pm1\}$.
Thus any difference of two signed components is
$A+B\cos t+C\sin t$.  If this vanishes near zero, its value and its
first two derivatives at zero give $A=B=C=0$.
Equality therefore holds for all real $t$, including any zero-radius
points.  Conversely, equal signed components give the regular orbit
in Eq.~\eqref{bound:orbit-four}, with radius
$\sqrt3\,|m_d|$.  This proves the local-to-global step used in
Theorem~\ref{bound:theorem}.

\subsection{Background terms and coefficient comparison}
Write $N=2^n$, $c=2/N$, and $q=d+\delta_0$.
Then $q$ is Boolean with $q(0)=1$, and Eq.~\eqref{core:defect-state}
applies with $\Theta=tq$. Combining the local formula
\eqref{core:signal} with the Boolean backgrounds in
Eq.~\eqref{core:boolean-background} gives the components for
classes I and II of Table~\ref{core:pullback-table}:
\begin{equation}
 m_I(t)=\begin{cases}
 c\cos t-W_q(h)\sin t,&q(a)=1,\\
 c-[c+W_q(h)]\sin t,&q(a)=0,
 \end{cases}
 \label{bound:component-i}
\end{equation}
and
\begin{equation}
 \begin{split}
 m_{II}(t)={}&-\tfrac12W_{D_aq}(h)(1-\cos t)\\
 &+c\bigl\{[1-q(a)]\cos t+q(a)(1-\sin t)\bigr\}.
 \end{split}
 \label{bound:component-ii}
\end{equation}

For $n\ge3$, both classes occur.  The cosine coefficient in
Eq.~\eqref{bound:component-i} forces all class-I flip values to equal
some $s\in\{0,1\}$.  The sine coefficient in
Eq.~\eqref{bound:component-ii} forces all class-II flip values to
equal some $k\in\{0,1\}$.  Comparing the sine coefficients across
the classes gives $W_q(h)=c(s+k-1)$, while comparing the constants
gives $W_{D_aq}(h)=2c(s+k-1)$.  The cosine coefficients then also
agree.  Substitution proves both directions of
Eq.~\eqref{bound:conditions} and yields the four laws.
Their maximum absolute values prove the general upper bound.

\subsection{Analytic exclusions at \texorpdfstring{$n=2,3,4$}{n=2,3,4}}
Use binary integers with $z_1$ the most significant bit.
For $n\ge3$ and every class-I character, the Walsh condition is equivalently
\begin{equation}
 \sum_z q(z)(-1)^{h\cdot z}=1-s-k.
 \label{bound:linear-character}
\end{equation}
Indeed, $W_q(h)=-2N^{-1}\sum_zq(z)(-1)^{h\cdot z}$ for $h\ne0$.

For $n=2$, only class I occurs and its flip masks cover all three
nonzero strings.  Equality of the cosine coefficients forces
$q(z)=s$ for every $z\ne0$.  The choices $s=0$ and $s=1$
give $q=\delta_0$ and $q\equiv1$, respectively.
Their common components are $c$ and $c\cos t$, so neither exceeds
the base radius in absolute value.

For $n=3$, the class-I flip masks are $\{1,2,3,4,5\}$ and
the class-II flip masks are $\{6,7\}$.
Thus $q(0)=1$, the first set has value $s$, and the second set
has value $k$; there are no further unknowns.
For $h=1$, the left side of Eq.~\eqref{bound:linear-character}
is $1-s$, which gives $k=0$.
For $h=4$, it is $1+s-2k$, which gives $2s-k=0$.
Hence $s=k=0$ and $d=0$.

For $n=4$, the class-I flip masks are $\{1,2,3,8,9\}$ and
the class-II masks are $\{6,7,12,13\}$.  Put
\begin{equation}
 (A,B,C,D,E,F)=(q_4,q_5,q_{10},q_{11},q_{14},q_{15}),
 \label{bound:unknown-values}
\end{equation}
where $q_j=q(j)$.  The four class-I characters $h=1,3,7,15$
give, respectively,
\begin{equation}
 \begin{aligned}
 A-B+C-D+E-F+k&=0,\\
 A-B-C+D-E+F+k&=0,\\
 -A+B-C+D+E-F+k&=0,\\
 -A+B+C-D-E+F+k&=0.
 \end{aligned}
 \label{bound:four-boundary-equations}
\end{equation}
Adding these equations gives $4k=0$ and therefore $k=0$.
Equation~\eqref{bound:four-laws} now leaves only $m=c$ or
$m=c\cos t$.  This exclusion covers every Boolean function at $n=4$; no ansatz for its support was used.
It is sufficient for the radius bound, so a list of individual
admissible supports is unnecessary.

The low-dimensional upper bounds are attained at $t=0$.
For every $n\ge5$, the explicit unified generator satisfies the
admissibility conditions and gives $m(-\pi/2)=2c_n$.
These two observations complete the proof of
Eq.~\eqref{bound:sharp-max}, with a maximum rather than merely a
supremum.
